\section{Certification and reproducibility}
\label{sec:certification}

The classification established in the preceding sections uses several
logically distinct computations. This section records their roles and
the checks connecting them.

The purpose of the certification layer is not to replace proof by
software output. Rather, it makes every finite claim reproducible from
the exact labeled graph and exposes the boundaries between construction,
enumeration, identification, and verification.

\subsection{Native graph certificate}

The native graph certificate records:

\begin{itemize}
\item the ordered list of sixty vertices;
\item the complete list of $120$ undirected edges;
\item the quartic degree check;
\item connectedness;
\item the canonical deck involution $a$;
\item the thirty $a$-fibers.
\end{itemize}

These data determine the exact labeled graph used throughout the paper.

\subsection{Quotient certificate}

The quotient certificate records the projection
\[
\pi\colon G60\longrightarrow G30,
\]
including:

\begin{itemize}
\item the thirty quotient vertices;
\item the assignment of each native vertex to its $a$-fiber;
\item the quotient edge set;
\item the exact quotient automorphism permutations;
\item the identification
\[
|\operatorname{Aut}(G30)|=240.
\]
\end{itemize}

The lift audit checks every quotient automorphism against the complete
native edge set.

Its result is:

\[
240\text{ quotient automorphisms},
\]
\[
2\text{ native lifts per quotient automorphism},
\]
and therefore
\[
480\text{ lifted native automorphisms}.
\]

\subsection{Independent native enumeration}

The native automorphism enumeration is performed directly on the
sixty-vertex graph.

It does not assume:

\begin{itemize}
\item preservation of the $a$-fiber partition;
\item centrality of the deck involution;
\item completeness of the lifted subgroup;
\item any proposed abstract group structure.
\end{itemize}

The independent enumeration returns exactly
\[
480
\]
native automorphisms.

The enumerated permutation set is compared directly with the lifted
permutation set. The comparison gives
\[
\text{native only}=0,
\qquad
\text{lifted only}=0.
\]

Thus the two sets agree exactly.

\subsection{Fiber and deck checks}

For every native automorphism $g$ and every native vertex $v$, the
certificate checks that
\[
\{g(v),g(a(v))\}
\]
is an $a$-fiber.

It also checks the permutation identity
\[
ga=ag
\]
for every
\[
g\in\operatorname{Aut}(G60).
\]

These tests establish that every native automorphism preserves the
canonical quotient and that the deck involution generates the center.

\subsection{Group-invariant certificate}

The complete native group is used to compute:

\begin{itemize}
\item the center;
\item the derived subgroup;
\item the second derived subgroup;
\item the abelianization;
\item the element-order profile;
\item the three index-two subgroups.
\end{itemize}

The certified values are
\[
|Z(A)|=2,
\]
\[
|A'|=120,
\qquad
A'\cong A_5\times C_2,
\]
\[
A''\cong A_5,
\]
and
\[
A/A'\cong C_2\times C_2.
\]

The element-order profile is
\[
\begin{array}{c|rrrrrrrr}
\text{order}
&1&2&3&4&5&6&10&12\\
\hline
\text{count}
&1&83&20&140&24&100&72&40.
\end{array}
\]

\subsection{Fiber-product comparison}

Reference groups are generated for the possible dihedral quotient
characters.

The cyclic-kernel candidate is rejected.

The Klein-four-kernel candidate agrees with the native group in:

\begin{itemize}
\item order;
\item center;
\item derived series;
\item abelianization;
\item element-order profile;
\item index-two subgroup structure.
\end{itemize}

This determines the reference model
\[
S_5\times_{C_2}D_8
\]
used in the explicit isomorphism certificate.

\subsection{Explicit isomorphism certificate}

The explicit certificate records:

\begin{itemize}
\item the selected native generators;
\item the corresponding fiber-product generators;
\item the complete $480$-element bijection;
\item the native and model encodings of every mapped element;
\item the multiplication verification count;
\item the canonical mapping digest.
\end{itemize}

The verified multiplication count is
\[
480^2=230400.
\]

The number of failed multiplication checks is
\[
0.
\]

The canonical mapping digest is
\[
\texttt{fdd62c4eb77eb3f553987ac415c1f6d00b434c56df4f2ae88fef90b3b6f6515d}.
\]

\subsection{Logical dependency order}

The classification follows the dependency chain
\[
\text{native graph}
\longrightarrow
\text{quotient and lifts}
\longrightarrow
\text{independent enumeration}
\longrightarrow
\text{group invariants}
\longrightarrow
\text{fiber-product model}
\longrightarrow
\text{explicit isomorphism}.
\]

No later identification is used to justify an earlier enumeration.

In particular:

\begin{itemize}
\item the order $480$ is established before the abstract group is chosen;
\item fiber preservation is tested on the independently enumerated group;
\item the fiber-product model is compared only after native invariants are known;
\item the explicit isomorphism is verified only after both groups have
been independently constructed.
\end{itemize}

This separation is the principal reproducibility boundary of the paper.
